% =====================================================================
%  src/preamble.tex  —  导言区
%  由 Lead 维护；队友只读。若缺宏，在文末「追加区」补 \newcommand，
%  并在回报中列出新增宏名（不得改动上面的已有定义）。
% =====================================================================
\usepackage[a4paper,margin=2.6cm,top=3.0cm,bottom=2.8cm]{geometry}
\usepackage{amsmath,amssymb,amsthm,mathtools,amscd}
\usepackage{mathrsfs}
\usepackage{graphicx}
\usepackage{array,booktabs,longtable}
\usepackage{caption}
\usepackage{enumitem}
\usepackage{xcolor}
\usepackage{fancyhdr}
\usepackage{tikz-cd}   % 交换图（原文若有）

% ---------- 字体（Windows，显式路径，不依赖 fontconfig） ----------
\setCJKmainfont[Path=C:/Windows/Fonts/,BoldFont=simhei.ttf,ItalicFont=simkai.ttf]{simsun.ttc}
\setCJKsansfont[Path=C:/Windows/Fonts/]{simhei.ttf}
\setCJKmonofont[Path=C:/Windows/Fonts/]{simfang.ttf}
\setCJKfamilyfont{zhkai}[Path=C:/Windows/Fonts/]{simkai.ttf}
\setCJKfamilyfont{zhhei}[Path=C:/Windows/Fonts/]{simhei.ttf}
\setCJKfamilyfont{zhfang}[Path=C:/Windows/Fonts/]{simfang.ttf}
\providecommand{\kaishu}{\CJKfamily{zhkai}}
\providecommand{\heiti}{\CJKfamily{zhhei}}
\providecommand{\fangsong}{\CJKfamily{zhfang}}

% ---------- 数学宏（只用这些，禁止在块文件里自造） ----------
\newcommand{\dd}{\mathrm{d}}
\newcommand{\ee}{\mathrm{e}}
\newcommand{\ii}{\mathrm{i}}
\DeclareMathOperator{\Tr}{Tr}
\DeclareMathOperator{\Aut}{Aut}
\DeclareMathOperator{\id}{id}
\DeclareMathOperator{\ind}{ind}
\DeclareMathOperator{\Ham}{Ham}
\DeclareMathOperator{\Hilb}{Hilb}
\DeclareMathOperator{\Vect}{Vect}
\DeclareMathOperator{\Alg}{Alg}
\DeclareMathOperator{\sHam}{sHam}
\DeclareMathOperator{\SQM}{SQM}
\DeclareMathOperator{\QM}{QM}
\DeclareMathOperator{\QFT}{QFT}
\DeclareMathOperator{\qft}{qft}
\DeclareMathOperator{\sqft}{sqft}
\DeclareMathOperator{\EBord}{EBord}
\DeclareMathOperator{\Bord}{Bord}
\DeclareMathOperator{\Rep}{Rep}
\DeclareMathOperator{\Th}{Th}
\DeclareMathOperator{\BPU}{BPU}
\DeclareMathOperator{\MP}{MP}
\DeclareMathOperator{\Spec}{Spec}
\newcommand{\TMF}{\mathrm{TMF}}
\newcommand{\tmf}{\mathrm{tmf}}
\newcommand{\KO}{\mathrm{KO}}
\newcommand{\KR}{\mathrm{KR}}
\newcommand{\MU}{\mathrm{MU}}
\newcommand{\MSpin}{\mathrm{MSpin}}
\newcommand{\MString}{\mathrm{MString}}
\newcommand{\R}{\mathbb{R}}
\newcommand{\C}{\mathbb{C}}
\newcommand{\Z}{\mathbb{Z}}
\newcommand{\Q}{\mathbb{Q}}
\newcommand{\N}{\mathbb{N}}
\newcommand{\PP}{\mathbb{P}}
\newcommand{\HH}{\mathbb{H}}
\newcommand{\cA}{\mathcal{A}}
\newcommand{\cB}{\mathcal{B}}
\newcommand{\cC}{\mathcal{C}}
\newcommand{\cD}{\mathcal{D}}
\newcommand{\cE}{\mathcal{E}}
\newcommand{\cF}{\mathcal{F}}
\newcommand{\cG}{\mathcal{G}}
\newcommand{\cH}{\mathcal{H}}
\newcommand{\cL}{\mathcal{L}}
\newcommand{\cM}{\mathcal{M}}
\newcommand{\cN}{\mathcal{N}}
\newcommand{\cO}{\mathcal{O}}
\newcommand{\cP}{\mathcal{P}}
\newcommand{\cS}{\mathcal{S}}
\newcommand{\cU}{\mathcal{U}}
\newcommand{\cV}{\mathcal{V}}
\newcommand{\cX}{\mathcal{X}}
\newcommand{\Ho}{H^{\bullet}}
\newcommand{\eps}{\varepsilon}
\newcommand{\ol}[1]{\overline{#1}}
\newcommand{\figomit}[2]{\par\smallskip\noindent\begingroup\small\kaishu\color{gray}〔原文图 #1 从略：#2〕\endgroup\par\smallskip}

% ---------- 定理类环境（编号一律手工给：\begin{definition}[2.1]） ----------
\theoremstyle{definition}
\newtheorem*{definition}{定义}
\newtheorem*{example}{例}
\newtheorem*{exercise}{习题}
\newtheorem*{conjecture}{猜想}
\newtheorem*{hypothesis}{假设}
\newtheorem*{question}{问题}
% 中文正文不用 plain 斜体：定理/命题/推论/引理一律直立（definition 风格）
\theoremstyle{definition}
\newtheorem*{theorem}{定理}
\newtheorem*{proposition}{命题}
\newtheorem*{corollary}{推论}
\newtheorem*{lemma}{引理}
\theoremstyle{remark}
\newtheorem*{remark}{注}

% 让 \begin{definition}[2.1] 印成「定义 2.1」
\makeatletter
\def\thmhead@plain#1#2#3{%
  \thmname{#1}\thmnumber{\@ifnotempty{#2}{ #2}}%
  \thmnote{\@ifnotempty{#3}{\hspace{0.35em}#3}}}
\let\thmhead\thmhead@plain
\makeatother

% ---------- 版式 ----------
\setlength{\parindent}{2em}
\linespread{1.18}
% 中文与公式混排行内偶有难以断行的长串：给足紧急伸长量，避免 Overfull
\setlength{\emergencystretch}{2em}
\setlist{nosep,leftmargin=2.2em}
\renewcommand{\arraystretch}{1.15}

\pagestyle{fancy}
\fancyhf{}
\fancyhead[L]{\small\kaishu 椭圆上同调与量子场论}
\fancyhead[R]{\small\thepage}
\renewcommand{\headrulewidth}{0.4pt}
\fancypagestyle{plain}{\fancyhf{}\fancyfoot[C]{\small\thepage}\renewcommand{\headrulewidth}{0pt}}

\ctexset{
  section = {
    format = \Large\bfseries\heiti,
    beforeskip = 2.2ex plus 1ex minus .2ex,
    afterskip  = 1.4ex plus .2ex,
    name = {},
    number = \arabic{section},
  },
  subsection = {
    format = \large\bfseries\heiti,
    beforeskip = 1.8ex plus .6ex minus .2ex,
    afterskip  = 1.0ex plus .2ex,
  },
}

% 超链接（必须放最后）
\usepackage[hidelinks,bookmarksnumbered=false,bookmarksopen=true,
            pdftitle={椭圆上同调与量子场论（中文译注本）},
            pdfauthor={Daniel Berwick-Evans}]{hyperref}
\pdfstringdefDisableCommands{\def\ol#1{#1}}

% ===================== 追加区（队友可 append \newcommand） =====================
